% !TeX root = ../papers/nc-p3.tex
\section{二対の非単位面値と三次の同時変形障害}
\label{sec:quantum-double-pair}
面値を $(c,c,d,d)$, $c\ne d$, $c,d\ne1$ とする。$\gamma=cd\in\{1,-1\}$ である。
両対を区別して高い側の文字を $x,y$、他方を $z,w$ と呼び、後の対では純平方の文字を $z$ に取る。
この選び方を全て許す。$a^2=d$ の一つを固定すれば、量子背景の三角積判定により
\[
 yx=xy,\quad zx=a xz,\quad zy=a^{-1}yz,\quad
 wx=a^{-1}xw,\quad wy=a yw,\quad wz=\gamma zw
\]
を使える。
\begin{theorem}\label{thm:quantum-double-pair}
この量子点へ退化する全窓は、下記の通常の次数付きOre表示で尽きる。
$yx=xy+D z^2$ とし、$zx,zy$ の二関係は背景の式を用いる。
$w$ による拡大は $w v=\sigma(v)w+\delta(v)$,
$(\sigma(x),\sigma(y),\sigma(z))=(a^{-1}x,ay,\gamma z)$ である。
\[
 \delta(z)=Ax^2+Bxy+Cy^2,\qquad
 \delta(x)=u xz+v yz,\quad \delta(y)=s xz+t yz,
\]
\[
 s=\frac{DA(1+\gamma a^2)}{a-a^{-1}},\qquad
 v=\frac{DC(1+\gamma a^{-2})}{a-a^{-1}} .
\]
$\gamma=1$ では条件は $BD=0$、$u=t=0$ とする。
$\gamma=-1$ では $A,B,C,D$ は全て自由で、$u=DB/(2a)$, $t=-DB/(2a)$ とする。
全類で最小周期は1であり、特殊値も標準Koszul左右AS正則である。
\end{theorem}
\begin{proof}[四欄の全被覆]
非零な横断接指標は前の面値計算により二つの三つ組だけであり、各々一次元以下である。
その三つ組は $x,y$ の二次単項式の次数から $e_z+e_w$ を引いたもの、
および $z,w$ の二次単項式の次数から $e_x+e_y$ を引いたものである。
それぞれを一組だけ使用した全三方向は一関係の単独書換えで実現し、各次元は1である。
実際、$wz$ の補正だけなら $\delta(z)$ は $x,y$ の任意の二次式でよい。
$yx$ の補正だけなら二対を交換した表示を使える。
両対の余重みの平均を比較し、$x,y$ 側を高く取ると、低い側の三つ組で正となり得るのは
余重みが最大の文字の純平方だけである。平均が等しい場合にも同じ包含が成立する。
従って流入点の横断接座標の非零台は $A,B,C,D$ の四欄に入る。

この場合、余重みは改めて $(1,1,2,-2)$ にできる。実際、切片の接空間への
トーラス同変線型射影 $\pi$ を接空間上の恒等と取ると、完備局所環はその座標で位相的に生成される。
元の流入弧で、重み零座標は背景値に等しく、他の非零座標は四欄だけである。
各テンソル座標をこの弧上の単項式で表すと、元の余重みが正なので各係数への寄与は有限で、
その多重次数は使った四欄の多重次数の非負和に限られる。
四欄はいずれも新しい余重みで2であり、従って新しい余指標でも同じ点が流入する。
同じ四値を持つ切片点は元の完備座標での弧の一致により一意である。

背景を通常の次数付き表示から標準四窓表示へ移す位置別対角変換は固定してよい。
この同定で、各欄の一次変形は上の通常関係の対応する係数で測る。
一つずつの補正には文字変更による一次境界がない（負添字が二つ）ので、
四係数と四接座標の対応は独立な非零倍率だけである。
通常のOre表示から得る族を切片上に戻したときも、この一次対応がそのまま
余重み2の四座標を与えることを確認しておく。
各次数で順序語が基底の族では関係と双対の最高次直線から標準四窓を基点の近くで正則に取れる。
最高次の倍率は体積の背景係数、四移動の基底は双対積で規格化でき、対角文字作用と可換する。
周囲の軌道接のLie補空間を対角同変に選ぶと、
その形式指数による基底変更とアフィン横断面の積から周囲への写像は一次で同型である。
従って形式逆を用いて切片点と基底変更を同変に取れる。
入力が全て余重み2なので余重み2の出力には線型項しか入らない。
所定の四値で得る二つの流入弧がこの形式同型で対応すれば、その一般点は同型である。
定数の二窓に戻した同型方程式は $k((u))$ で解を持つため、代数閉体 $k$ でも解を持つ。
以上により全四値のOre族を構成できた場合、または必要な障害を除いて構成できた場合に、
それ以上の窓は残らない。
\end{proof}

\begin{proof}[重なりと障害の必要性]
三文字部分は $z$ を先に置いた反復Ore拡大である。例えば $zx=a xz$ の量子平面へ
$y$ を加え、$yx=xy+D z^2$, $yz=a zy$ を用いれば基底ができる。
$\sigma$ はこの三関係を保つ。
途中では $\delta(z)=Ax^2+Bxy+Cy^2+H z^2$ を許す。
$\delta$ を順に $zx-a xz$, $zy-a^{-1}yz$, $yx-xy-Dz^2$ に適用して順序語の係数を取ると
\begin{align*}
 DB+H(a^2-1)+a(\gamma-1)u&=0,& DC(1+a^{-2})+(\gamma/a-a)v&=0,\\
 -DA(1+a^2)+(\gamma a-a^{-1})s&=0,&-DBa^{-2}+H(a^{-2}-1)+(\gamma-1)t/a&=0,\\
 (a-a^{-1})s&=DA(1+\gamma a^2),&(a-a^{-1})v&=DC(1+\gamma a^{-2}),\\
 (a-a^{-1})(t+u)&=DB(1+\gamma),&D\{a(t+u)-(1+\gamma)H\}&=0 .
\end{align*}
$\gamma=-1$ では定理の値と $H=0$ で全式が消える。
$\gamma=1$ では $H=-DB/(a^2-1)$ と $t+u=2DB/(a-a^{-1})$ を最後に代入すると
$2D^2B(a^2+1)/(a^2-1)=0$ が必要である。$a^2=1$ は非単位条件、
$a^2=-1$ はこの符号の場合の $c\ne d$ に反する。従って $DB=0$ で、
この条件下では定理の $H=u=t=0$ が全式を満たす。

$\gamma=1$ において最後の障害がOre表示の仮定を外しても必要であることを示す。
この場合普通表示の行積は1であり、背景体積は符号付き巡回不変である。
一関係だけの各一次変形も位置モード1である。実際、二つの相補文字の内部の交換係数は1で、
他の二文字との交換係数の積も1であるから、前節の $1-\eta\chi_j$ で $\eta=1$ となる。
巡回同変な切片と接射影を使えば、巡回した点も同じ四座標を持つ。
流入弧の一意性によりその点自体が巡回不変で、六関係を全位置で同じ表示に取れる。
順序を $x,y,z,w$ とし、逆順対を左辺とする六関係を行消去で規格化する。
必要なPlücker座標は背景で非零であり、多重次数を持つので流入点でも非零である。
各補正は四欄の重み2の非負和を持つ。二文字の余重みを並べれば、
$yx$ に加え得るのは $z^2$、$zx,zy$ には何もなく、$wx,wy$ には $xz,yz$、
$wz$ には $x^2,xy,y^2,z^2$ だけである。従ってその六関係は $H$ を許した上記の一般式に尽きる。
三文字の順序語は独立である。量子極限への各経路は命題\ref{prop:unrestricted-good-reduction}
により自由で、特殊点の順序語基底が持ち上がるためである。
従って $w$ と各三文字関係の重なりから八式を必ず満たす。
これで必要性を得た。逆に定理の係数では実際に捩れ導分を定めたのでOre拡大の
順序語が全次数で基底である。随伴次数を量子多項式環にする重み（または反復する
斉次Ore拡大の線型解消）から標準Koszul左右AS正則性が従い、表示は一位置周期である。
\end{proof}

\begin{corollary}
全特殊値では、二つの三つ組の非零台と、三項が揃った三つ組での
「両平方係数の積を混合係数の二乗で割った値」が全トーラス不変量である。
さらに $\gamma=1$ なら二組を交換する四巡回とその冪、
$\gamma=-1$ なら両組内の同時反転で台を同一視する。
同じファイバー内の退化はこれらの置換と台の露出面への制限で尽くされる。
$\gamma=1$ の極大台は指定を忘れると二型で同型類次元は1と0、
$\gamma=-1$ でも二型でどちらも1である。
\end{corollary}
\begin{proof}
正の側の三指標は端の和が中央の二倍という関係だけを持ち、他方の一純平方は
その有理線型包の外にある。従って整数核の不変量は述べた一比だけである。
面の向き付き置換条件から偶置換は恒等と $(xy)(zw)$ だけで、
奇置換が許されるのは $d=c^{-1}$ のときだけである。そのときの $(xzyw)$ は上記の
交換係数も保つ。各置換による左辺の並べ替えの共通倍率は比から消え、
混合単項式の並べ替えの倍率は $1$ または $-1$ で二乗が1である。
より高い余重みの項は四欄の値に寄与しない。切片ではこれらの座標と全余重み零生成元で
流入弧が決まるので、軌道閉包への面の条件もトーラスと同じである。
正符号では三項だけの台と、二平方および反対側の一平方の台が極大で、
後者の三指標は独立である。負符号では三項と反対側の一平方を持つ台が極大で、
三項をどちら側に置くかは同一視されない。
\end{proof}
